% Plantilla adaptada desde investigacion-aplicada-a-la-contaduria/reporte-foro-S1.tex
% Compilar con: latexmk -pdf -pdflatex="lualatex %O %S"
\makeatletter
% Añadir rutas de búsqueda para imágenes y plantillas (relativas a este archivo)
\graphicspath{{../investigacion-aplicada-a-la-contaduria/plantillas/}{../investigacion-aplicada-a-la-contaduria/}{./assets-actividades/}{./}}
\makeatother
% Base reutilizable; no compilar directamente. LuaLaTeX + Biber.
% Compilar desde la carpeta de la asignatura mediante compilar-pendientes.sh.
\documentclass[12pt,letterpaper,oneside]{article}
\usepackage[margin=2.54cm]{geometry}
\usepackage{fontspec}
\setmainfont{Arial}
\setsansfont{Arial}
\usepackage[spanish,es-tabla,es-nodecimaldot]{babel}
\usepackage{csquotes,setspace,ragged2e}
\usepackage{graphicx}
\usepackage[table]{xcolor}
\usepackage{booktabs,array,longtable,tabularx}
\usepackage{enumitem,titlesec}
\usepackage{tikz}
\usetikzlibrary{arrows.meta,positioning,calc}
\usepackage{pdflscape}
\usepackage[most]{tcolorbox}
\usepackage[backend=biber,style=apa,sortcites=true]{biblatex}
\DeclareLanguageMapping{spanish}{spanish-apa}
\addbibresource{fuentes-semanales.bib}
\usepackage{xurl}
\usepackage[hidelinks]{hyperref}
\usepackage{attachfile2}
\definecolor{uasBlue}{HTML}{123B6D}
\definecolor{uasGold}{HTML}{A77A16}
\definecolor{uasPaper}{HTML}{F2F6FA}
\definecolor{uasGreen}{HTML}{276B53}
\hypersetup{pdfauthor={Martín Jonathan De La Cruz Muñoz},
  pdfsubject={Investigación aplicada a la contaduría — UAS},hypertexnames=false}
\onehalfspacing
\setlength{\parindent}{1.27cm}
\setlength{\parskip}{.3em}
\setlength{\emergencystretch}{3em}
\setlength{\bibhang}{1.27cm}
\setlength{\bibitemsep}{.6em}
\raggedbottom
\titleformat{\section}{\normalsize\bfseries\color{uasBlue}}{\thesection.}{.5em}{}
\titleformat{\subsection}{\normalsize\bfseries}{\thesubsection}{.5em}{}
\setlist{leftmargin=*,itemsep=.3em,parsep=0pt}
\newcolumntype{Y}{>{\RaggedRight\arraybackslash}X}
\newcolumntype{L}[1]{>{\RaggedRight\arraybackslash}p{#1}}
\newcommand{\fechaRealizacion}{13 de septiembre de 2026}
\newcommand{\portada}[2]{%
  \hypersetup{pdftitle={#1}}%
  \begin{titlepage}
    \centering
    \includegraphics[height=2.2cm]{assets-actividades/escudo-uas.png}\hfill
    \includegraphics[width=6cm]{../fcauas-logo.jpg}\par
    \vspace{.8cm}
    {\normalsize\bfseries Universidad Autónoma de Sinaloa\par}
    Facultad de Contaduría y Administración\par
    Licenciatura en Contaduría Pública — Modalidad Virtual\par
    \vfill
    {\normalsize\bfseries\color{uasBlue}#1\par}
    \vspace{.6cm}#2\par
    \vspace{.6cm}{\normalsize Investigación aplicada a la contaduría\par}
    \vfill
    \begin{tabularx}{\textwidth}{@{}lY@{}}
      Estudiante: & Martín Jonathan De La Cruz Muñoz\\
      Matrícula: & 26031061\\
      Asesora: & Dra. Nadia Aileen Valdez Acosta\\
      Grupo: & 103 LCP\\
      Fecha de realización: & \fechaRealizacion\\
      Lugar: & Culiacán, Sinaloa, México
    \end{tabularx}
  \end{titlepage}
}
\newtcolorbox{forobox}[1]{enhanced,breakable,title={#1},
  colback=uasPaper,colframe=uasBlue,colbacktitle=uasBlue,coltitle=white,
  fonttitle=\bfseries,boxrule=.6pt,arc=1mm,
  left=3mm,right=3mm,top=2mm,bottom=2mm,
  before upper={\setlength{\parindent}{0pt}\setlength{\parskip}{.5em}}}
\attachfilesetup{color={0.071 0.231 0.427},print=false}
\newcommand{\foroTexto}[1]{\noindent\textattachfile[
  description={Texto de la intervención para revisar y copiar},mimetype=text/plain]{#1}{%
  \colorbox{uasBlue}{\textcolor{white}{\small Abrir texto adjunto}}}\par}
\newcommand{\fuenteCaja}[1]{\par\begingroup\small\fullcite{#1}\par\endgroup}
\newcommand{\referencias}{\clearpage\printbibliography[title={Referencias}]}
\tikzset{
  raiz/.style={rounded corners,draw=uasBlue,fill=uasBlue,text=white,
    align=center,text width=8cm,inner sep=8pt,font=\bfseries},
  concepto/.style={rounded corners,draw=uasBlue,fill=uasPaper,
    align=center,text width=5.6cm,inner sep=7pt,font=\normalsize},
  detalle/.style={rounded corners,draw=uasGreen,fill=white,
    align=left,text width=5.6cm,inner sep=7pt,font=\normalsize},
  enlace/.style={-{Stealth[length=2mm]},thick,draw=uasBlue},
  nexo/.style={fill=white,inner sep=2pt,font=\small,align=center}
}
\begin{document}
\portada{Reporte de Lectura — Introducción al Derecho}{Reporte de Actividad 1}

\section{Antecedentes Históricos del Derecho}
El Derecho surge históricamente como respuesta a la necesidad de regular la convivencia en sociedades crecientes. Desde normas consuetudinarias y las leyes escritas de civilizaciones antiguas, se observa la progresiva institucionalización de reglas para dirimir conflictos y organizar la vida colectiva. Estos antecedentes muestran la transición de la costumbre a sistemas jurídicos complejos y codificados.

\section{Acepciones de la Palabra Derecho}
El término Derecho tiene varias acepciones: como norma (conjunto de reglas), como ciencia (disciplina que estudia las normas) y como facultad (derechos subjetivos de las personas). Cada acepción refleja una dimensión distinta: normativa, teórica y práctica, y juntas permiten entender cómo se aplica el Derecho en la realidad social.

\section{Fuentes del Derecho}
Las fuentes del Derecho incluyen la ley, la costumbre, los principios generales del Derecho y la jurisprudencia. La ley escrita proporciona seguridad y previsibilidad; la costumbre regula prácticas sociales cuando la ley es silenciosa; la jurisprudencia interpreta y unifica criterios; y los principios orientan soluciones cuando las reglas concretas son insuficientes.

\section{Normas Jurídicas}
Las normas jurídicas son mandatos de conducta obligatorios y sancionados por el Estado; se caracterizan por su generalidad, abstracción y coercibilidad. Su función es ordenar comportamientos, establecer deberes y atribuir derechos, y su incumplimiento puede dar lugar a sanciones o reparaciones.

\section{El Estado}
El Estado es la estructura política que crea y aplica el Derecho; posee el monopolio legítimo de la fuerza y las instituciones necesarias para promulgar leyes, administrar justicia y garantizar seguridad. El vínculo entre Estado y Derecho es esencial: sin una autoridad que haga cumplir las normas, la efectividad de las reglas sería limitada.

\section{Conclusión}
El Derecho es indispensable en la vida diaria porque organiza la convivencia, protege derechos, brinda seguridad jurídica y ofrece mecanismos de resolución de conflictos. Comprender sus antecedentes, fuentes y la naturaleza de las normas permite aplicar principios jurídicos con mayor claridad y responsabilidad.

\referencias
\end{document}
